% GENERATED FILE. Edit canonical lessons, then run npm run generate:books.
% canonical-source-hash: fnv1a64:2a3618ec36296e56
% canonical-lessons: TE-C143-errati, TE-C143-egiru, TE-C143-parigettu, TE-C143-nidrapo, TE-C143-lecu

\chapter{Red, To fly, To run, To sleep, To get up}
\label{ch:te-red-to-fly-to-run-to-sleep-to-get-up}
\hlchaptermodality{143}

\begin{chapteropening}
\textbf{By the end of this chapter:} \emph{I can say red, to fly, to run, to sleep and to get up.}

Complete the last lesson of chapter 143: I can say red, to fly, to run, to sleep and to get up.
\end{chapteropening}

\section[\texorpdfstring{\te{ఎర్రటి} (erraṭi)}{ఎర్రటి (erraṭi)}]{\te{ఎర్రటి} (erraṭi) — red, reddish}

\label{lesson:TE-C143-errati}



\begin{quote}
Before the new one: say the Telugu for black, then the Telugu for white.
\end{quote}

\subsection*{You'll want to know: \te{ఎర్రటి}}
\textbf{\te{ఎర్రటి}} — \emph{erraṭi} — \textquotedblleft{}red, reddish\textquotedblright{}.

\begin{grammarlens}[title={What you've built}]
One more word for this chapter.
\end{grammarlens}

\subsection*{Your turn}
\begin{itemize}
\raggedright
  \item \emph{Say it:} \emph{erraṭi}
  \item \emph{Say it:} \emph{erraṭi}, once more
  \item \emph{Say it:} say \emph{erraṭi}
  \item \emph{Recall:} say the Telugu for black, then the Telugu for white, then say \emph{erraṭi} again
\end{itemize}

\begin{tcolorbox}[breakable,colback=teal!4,colframe=teal!35!black,title={Before you move on}]
What does \te{ఎర్రటి} mean? (\textquotedblleft{}Red, reddish\textquotedblright{}.) Say it once more.
\end{tcolorbox}
\section[\texorpdfstring{\te{ఎగురు} (egiru)}{ఎగురు (egiru)}]{\te{ఎగురు} (egiru) — to fly}

\label{lesson:TE-C143-egiru}



\begin{quote}
Before the new one: say the Telugu for white, then the Telugu for red.
\end{quote}

\subsection*{You'll want to know: \te{ఎగురు}}
\textbf{\te{ఎగురు}} — \emph{egiru} — \textquotedblleft{}to fly\textquotedblright{}.

\begin{grammarlens}[title={What you've built}]
One more word for this chapter.
\end{grammarlens}

\subsection*{Your turn}
\begin{itemize}
\raggedright
  \item \emph{Say it:} \emph{egiru}
  \item \emph{Say it:} \emph{egiru}, once more
  \item \emph{Say it:} say \emph{egiru}
  \item \emph{Recall:} say the Telugu for white, then the Telugu for red, then say \emph{egiru} again
\end{itemize}

\begin{tcolorbox}[breakable,colback=teal!4,colframe=teal!35!black,title={Before you move on}]
What does \te{ఎగురు} mean? (\textquotedblleft{}To fly\textquotedblright{}.) Say it once more.
\end{tcolorbox}
\section[\texorpdfstring{\te{పరిగెత్తు} (parigettu)}{పరిగెత్తు (parigettu)}]{\te{పరిగెత్తు} (parigettu) — to run}

\label{lesson:TE-C143-parigettu}



\begin{quote}
Before the new one: say the Telugu for red, then the Telugu for to fly.
\end{quote}

\subsection*{You'll want to know: \te{పరిగెత్తు}}
\textbf{\te{పరిగెత్తు}} — \emph{parigettu} — \textquotedblleft{}to run\textquotedblright{}.

\begin{grammarlens}[title={What you've built}]
One more word for this chapter.
\end{grammarlens}

\subsection*{Your turn}
\begin{itemize}
\raggedright
  \item \emph{Say it:} \emph{parigettu}
  \item \emph{Say it:} \emph{parigettu}, once more
  \item \emph{Say it:} say \emph{parigettu}
  \item \emph{Recall:} say the Telugu for red, then the Telugu for to fly, then say \emph{parigettu} again
\end{itemize}

\begin{tcolorbox}[breakable,colback=teal!4,colframe=teal!35!black,title={Before you move on}]
What does \te{పరిగెత్తు} mean? (\textquotedblleft{}To run\textquotedblright{}.) Say it once more.
\end{tcolorbox}
\section[\texorpdfstring{\te{నిద్రపో} (nidrapō)}{నిద్రపో (nidrapō)}]{\te{నిద్రపో} (nidrapō) — to sleep}

\label{lesson:TE-C143-nidrapo}



\begin{quote}
Before the new one: say the Telugu for to fly, then the Telugu for to run.
\end{quote}

\subsection*{You'll want to know: \te{నిద్రపో}}
\textbf{\te{నిద్రపో}} — \emph{nidrapō} — \textquotedblleft{}to sleep\textquotedblright{}.

\begin{grammarlens}[title={What you've built}]
One more word for this chapter.
\end{grammarlens}

\subsection*{Your turn}
\begin{itemize}
\raggedright
  \item \emph{Say it:} \emph{nidrapō}
  \item \emph{Say it:} \emph{nidrapō}, once more
  \item \emph{Say it:} say \emph{nidrapō}
  \item \emph{Recall:} say the Telugu for to fly, then the Telugu for to run, then say \emph{nidrapō} again
\end{itemize}

\begin{tcolorbox}[breakable,colback=teal!4,colframe=teal!35!black,title={Before you move on}]
What does \te{నిద్రపో} mean? (\textquotedblleft{}To sleep\textquotedblright{}.) Say it once more.
\end{tcolorbox}
\section[\texorpdfstring{\te{లేచు} (lēcu)}{లేచు (lēcu)}]{\te{లేచు} (lēcu) — to get up}

\label{lesson:TE-C143-lecu}



\begin{quote}
Before the new one: say the Telugu for to run, then the Telugu for to sleep.

Say eight commands, one breath each: \textbf{\te{రా}}, \textbf{\te{ఇవ్వు}}, \textbf{\te{కూర్చో}}, \textbf{\te{నిల్చో}}, \textbf{\te{ఆలకించు}}, \textbf{\te{చూడు}}, \textbf{\te{మాట్లాడు}}, \textbf{\te{రాయి}}.
\end{quote}

\subsection*{You'll want to know: \te{లేచు}}
\textbf{\te{లేచు}} — \emph{lēcu} — \textquotedblleft{}to get up\textquotedblright{}.

\begin{grammarlens}[title={What you've built}]
One more word for this chapter.
\end{grammarlens}

\subsection*{Your turn}
\begin{itemize}
\raggedright
  \item \emph{Say it:} \emph{lēcu}
  \item \emph{Say it:} \emph{lēcu}, once more
  \item \emph{Say it:} say \emph{lēcu}
  \item \emph{Recall:} say the Telugu for to run, then the Telugu for to sleep, then say \emph{lēcu} again
\end{itemize}

\begin{tcolorbox}[breakable,colback=teal!4,colframe=teal!35!black,title={Before you move on}]
What does \te{లేచు} mean? (\textquotedblleft{}To get up\textquotedblright{}.) Say it once more.
\end{tcolorbox}
